\section{Thermodynamic row kernels}\label{sec:thermodynamic}

The finite formulas of Section~\ref{sec:finite} reduce a row word to
sums over finitely many selected Bethe roots.  To pass to a plane
kernel we must control those sums even when several selected roots
approach the sparse ends of the root configuration.  We first
construct the real roots, then prove the inverse and determinant
estimates that justify this passage.  All constants in this Section
may depend on the fixed parameter $0<\lambda<\pi$.  They need not
remain bounded as $\lambda\downarrow0$ or $\lambda\uparrow\pi$.

\subsection{The three limiting procedures}

We use the notation $s,K,r,G,\xi$ of Section~\ref{sec:finite}, and put
\[
 h=\int_{\R}|K(u)|\dd u=|1-2\eta|<1,\qquad
 R_0=K*(I+K*)^{-1}.
\]
Here $K*$ denotes convolution on the line; the last expression means
the convolution kernel of the indicated operator.  Since $h<1$, its
Neumann series converges on bounded functions.  There is also a
number $a>0$ such that the same series converges with the weight
$e^{a|u|}$: choose $a$ so that
$\int e^{a|u|}|K(u)|\dd u<1$.  In particular $R_0$ and each of its
fixed derivatives are exponentially integrable.

The following limits have different purposes.
\begin{enumerate}
\item In the \emph{homogeneous plane limit}, all $b_j=0$ and $L$
 tends to infinity through even integers.  If $\lambda>\pi/2$ we
 use $n=L/2$.  If $0<\lambda\le\pi/2$, we use
 $n=L/2-\lfloor L^\epsilon\rfloor$, where $\epsilon>0$ is
 sufficiently small and fixed.  Its permitted size is specified in
 the proof of Lemma~\ref{lem:row-normalization}.
\item In the \emph{array limit}, which we use only for
 $\lambda>2\pi/3$, put
 \[
 b_j=\left(j-\frac{L-1}{2}\right)\frac d2,\qquad
 B=\frac{Ld}{2},\qquad n=L/2.
 \]
 First $B\to\infty$ through the allowed even widths, and then
 $d\downarrow0$.  A common real center $a_B$ of the row arguments
 may vary with $B$, provided
 $|a_B|\le(1-\gamma)B/2$ for a fixed $\gamma>0$.
 Differences of arguments range over a fixed compact set.
\item In the \emph{regular-filling limit}, also used subsequently
 only for $\lambda>2\pi/3$, the chain is homogeneous.  First
 $L\to\infty$ with $n/L\to\rho<1/2$, and then
 $\rho\uparrow1/2$.  We assert comparison with the plane state in
 this procedure for equal-time local matrices.
\end{enumerate}
For example, uniform array convergence to a function $F$ means
\begin{equation}\label{eq:array-limit-convention}
 \lim_{d\downarrow0}\limsup_{B\to\infty}
 \sup_{\substack{|a_B|\le(1-\gamma)B/2\\
                  \boldsymbol{x}\in\mathcal K}}
 |F_{d,B}(a_B+\boldsymbol{x})-F(\boldsymbol{x})|=0
\end{equation}
for every fixed compact parameter set $\mathcal K$ under discussion.
It is not necessary to assert a limit for every fixed $d$.
A uniform bound in this procedure means that there are $C,d_0>0$
such that for every $0<d<d_0$ there is $B_0(d)$ for which the bound
holds whenever $B\ge B_0(d)$.  The corresponding convention for
regular filling is $\lim_{\rho\uparrow1/2}\limsup_{L\to\infty}$.
Word lengths, matrices, derivative orders, and compact parameter
sets are fixed before these limits.  No estimate below is asserted
uniformly in an unbounded word length.

\subsection{Real roots, counting cells, and the sparse ends}

There are two constructions of roots. At exact half filling with
$K<0$, the last few roots remain sparse and need individual estimates.
At regular filling, and in the homogeneous near-half-filled sequence
used when $K\ge0$, even the terminal cells become dense. The next
lemma supplies both constructions and the quadrature measure that
will enter the finite-word formula.

\begin{lemma}[Construction and density of the roots]\label{lem:root-limits}
The zero-twist consecutive-root equations of
Lemma~\ref{lem:bethe-data} have distinct real solutions with the
following properties.
\begin{enumerate}
\item If $\lambda>\pi/2$, the solutions at half filling exist for
 every finite homogeneous or real-array chain.  Their counting
 derivative $\rho_L=\xi'$ is positive.  If
 \[
 p_L(u)=\sum_{j=0}^{L-1}r(u-b_j),
 \]
 then $\rho_L-p_L$ and their integrated discrepancy are uniformly
 bounded on the entire line.
\item In the homogeneous chain, let
 \[
 r_Q=(I+K_Q)^{-1}s,\qquad H_Q=(I+K_Q)^{-1}1,
 \]
 where $K_Q$ integrates over $[-Q,Q]$.  For all $Q\ge0$,
 \begin{equation}\label{eq:regular-density-comparison}
 c r(u)\le r_Q(u)\le C r(u),\qquad
 c\le H_Q(u)\le C\quad (|u|\le Q).
 \end{equation}
 The mass $\int_{-Q}^Q r_Q$ increases to $1/2$, and
 \[
 \frac12-\int_{-Q}^Q r_Q\asymp e^{-Q/\eta}.
 \]
 Choose $Q=Q_L$ by $\int_{-Q}^Q r_Q=n/L$.
 If $Q$ stays bounded with a positive limiting filling, or if
 $Lr(Q)$ tends to infinity at least as a fixed positive power of
 $L$, the roots exist for all sufficiently large $L$.  If $y_j^0$ are the mid-quantiles of $Lr_Q$, then
 $Lr_Q(y_j^0)|y_j-y_j^0|\le C$ uniformly in $j$, and
 \begin{equation}\label{eq:regular-counting-density}
 \xi'(u)=Lr_Q(u)+O(1)
 \end{equation}
 on the interval, including a bounded counting extension at its
 endpoints.
\item In all three limiting procedures the measures
 \[
 \nu_L=\sum_j\ell_j\delta_{y_j},\qquad
 \ell_j=\xi'(y_j)^{-1},
 \]
 have uniformly bounded mass on intervals of length one.  In a
 region whose slightly enlarged neighborhood has density at least
 $M$, their error against bounded Lipschitz tests on a fixed
 bounded interval is at most
 \begin{equation}\label{eq:local-quadrature}
 \frac{C}{M}\bigl(\|\phi\|_\infty+\Lip(\phi)\bigr).
 \end{equation}
 Thus $\nu_L$ converges locally to Lebesgue measure in the dense
 portions of the limits.
\end{enumerate}
At zero twist the homogeneous vectors constructed from these roots
are the sector Perron vectors.  The Gaudin matrices are nonsingular,
so each constructed finite root list has a real-analytic twist branch
near zero.
\end{lemma}

\begin{proof}
We first construct the exactly half-filled lists and control their
sparse ends. We then compare the finite-interval density with the
full-line density and use its quantiles to construct the dense lists.
Both constructions will give the same local quadrature estimate.

\smallskip\noindent\textbf{Exactly half-filled roots.}
Suppose $K<0$ and write $k=-K>0$.  The Jacobian of the
counting equations is
\[
 G_{ii}=\sum_j s(y_i-b_j)+\sum_{j\ne i}k(y_i-y_j),
 \qquad G_{ij}=-k(y_i-y_j)\quad(i\ne j).
\]
Its quadratic form is the sum of a positive diagonal form and
$\sum_{i<j}k(y_i-y_j)|u_i-u_j|^2$.  It is strictly positive,
so the equations are strictly convex gradient equations.
Start at $\lambda=\pi/2$, where $K=0$ and each counting level has
one real preimage, and continue in a compact parameter interval.
There can be no coincident roots because their prescribed counting
levels are distinct.

There can be no escaping roots either.  If $k_0$ roots tend to
$+\infty$, sum their equations.  The interaction primitives within
that set cancel by oddness.  Dividing by $k_0$, the limiting excess
of the left side over the prescribed levels is
\[
 (1-\eta)(L/2-n+k_0)>0.
\]
The same argument applies at $-\infty$.  All site parameters here
are finite in each continuation problem.  Nonsingularity and this
compactness continue the solution to the required parameter.

At half filling the total mass of $\rho_L=\xi'$ is $n$.
The two outer counting half-levels occur at infinity.  Consequently
the signed measure and its primitive
\[
 e_L=\sum_j\delta_{y_j}-\rho_L(u)\dd u,\qquad
 F_L(u)=e_L((-\infty,u])
\]
satisfy $|F_L|\le1/2$ and $F_L(\pm\infty)=0$.
The density equation and convolution inversion give
\begin{equation}\label{eq:counting-sawtooth}
 \rho_L-p_L=-R_0*e_L=-R_0'*F_L.
\end{equation}
Integrating once also bounds the difference of the tail integrals
by $\|R_0\|_1/2$.  This proves the first assertion.  Moreover
$|\rho_L'|\le C\rho_L$: both $s$ and $k$ are positive kernels
with bounded logarithmic derivatives.

\smallskip\noindent\textbf{The sparse terminal cells.}
The preceding discrepancy estimate identifies the dense bulk, but a
uniform inverse must also control the last few roots. We therefore
estimate their densities and gaps directly.  Number the roots from the right as
$x_1>x_2>\cdots$ and write
$T(u)=\int_u^\infty\rho_L(v)\dd v$.  Then
$T(x_j)=j-1/2$.  For either $q=s$ or $q=k$, put
$\bar q(u)=\int_u^\infty q(v)\dd v$.
The displayed hyperbolic formulas for the kernels imply
\begin{align}
 q(u)&\le C\bar q(u) &&(u\in\R),\notag\\
 c\bar q(u)&\le q(u),\qquad
 \bar q(u+t)\le C e^{-2t}\bar q(u)
 &&(u,t\ge0).
 \label{eq:kernel-tail-comparison}
\end{align}
At a root $x_j$ lying beyond every bare center, its own kernel and
the $j-1$ farther-out kernels contribute at most
$h/2+h(j-1)=hT(x_j)$ to this tail.  The remaining centers are
inward.  Their tail is at least $(1-h)T(x_j)$, so
\begin{equation}\label{eq:edge-index-density}
 c(j-\tfrac12)\le\rho_L(x_j)\le C(j-\tfrac12).
\end{equation}
This bounds each fixed-index reciprocal density away from both
zero and infinity.

Terminal gaps cannot become long.  If $x_{j+1}$ also lies beyond
the bare centers and $\Delta=x_j-x_{j+1}$, compare the inward
tail at $x_j$ with its tail at $x_{j+1}$ using
\eqref{eq:kernel-tail-comparison}.  It gives
\[
 (1-h)(j-\tfrac12)
 \le C e^{-2\Delta}(j+\tfrac12),
 \qquad
 \Delta\le\frac12\log\frac{3C}{1-h}.
\]
Conversely $1=\int_{x_{j+1}}^{x_j}\rho_L$ and the upper tail
comparison give $\Delta\ge c/(j+1)$.
Thus each fixed-index gap has positive finite subsequential limits.

The explicit tail of $p_L$, together with its bounded integrated
error, locates any sufficiently large fixed edge index at
$e_L^{\rm edge}+O(1)$, where
\begin{equation}\label{eq:nominal-edges}
 e_L^{\rm edge}=\eta\log L
 \quad\hbox{or}\quad
 e_{d,B}^{\rm edge}=B/2+\eta\log(1/d)
\end{equation}
in the homogeneous and array cases, respectively.  For example,
choose an index larger than the integrated error bound; the
$p_L$ tail is comparable to
$\exp(-(u-e_L^{\rm edge})/\eta)$ near the displayed edge.
That root lies beyond the bare centers in the limiting procedures.
The gap estimate then places the last root no more than a bounded
distance beyond that edge.  The left end is identical.
For each fixed density threshold $M$, the portions with density
below $M$ start within a distance $C_M$ of the nominal edges and
contain only a bounded number of root indices.  Both the number and
the distance may depend on $M$.

These statements have their asserted array uniformity.  In a
purported failure sequence we may take $d\downarrow0$ and $B$
arbitrarily large depending on $d$.  Every fixed edge index is then
eventually beyond all bare centers, and
\eqref{eq:edge-index-density} applies to every such index.
One should not instead take a fixed-$d$ edge limit and assert that
its reciprocal weights tend to zero inward: the subsequent
$d\downarrow0$ limit is essential for that assertion.

\smallskip\noindent\textbf{Finite-interval densities.}
We turn to the dense construction. Its reference configuration will
consist of quantiles of $Lr_Q$, so we first need positivity and a
relative comparison with $r$, uniformly as $Q$ grows. The inverse of
$I+K_Q$ exists on bounded functions because its perturbation has norm
at most $h<1$.
Comparison with the full-line equation gives
\begin{equation}\label{eq:finite-Q-correction}
 r_Q-r=(I+K_Q)^{-1}K_{\rm out}r,
\end{equation}
where $K_{\rm out}$ integrates over $\R\setminus[-Q,Q]$.
When $K<0$, write the right side as the negative of the exit kernel
$(I-k_Q)^{-1}k_{\rm out}$ applied to $r$.
That positive kernel has mass at most $h$: the constant $h$ is a
supersolution of its mass equation, since
$h-k_Qh-k_{\rm out}1=(1-h)k_Q1\ge0$.
For exterior $v$ and interior $u$, we have
$r(v)\le r(Q)\le r(u)$. Applying the exit-kernel mass bound to
this inequality gives $(1-h)r\le r_Q\le r$ on the interval.  The positive Neumann
series also gives $1\le H_Q\le(1-h)^{-1}$ in this case.

For $0<\lambda<\pi/2$, a direct bound on the alternating Neumann
series for $I+K_Q$ would lose positivity when $h$ is large.
Instead the full-line dressed kernel is positive.  To see this
explicitly, put
\[
 \vartheta=\frac{\pi\lambda}{\pi-\lambda},\qquad
 J_\lambda(u)=
 \frac{\sin\vartheta}
 {(\pi-\lambda)
   \bigl(\cosh(2\pi u/(\pi-\lambda))-\cos\vartheta\bigr)}.
\]
Here $0<\vartheta<\pi$, so $J_\lambda\ge0$.
Rescaling the bare-kernel transform in
\eqref{eq:kernel-transforms} gives
\[
 \widehat J_\lambda(p)
   =\frac{\sinh((\pi-2\lambda)p/2)}
          {\sinh((\pi-\lambda)p/2)}.
\]
Together with $\widehat r(p)=(2\cosh(\lambda p/2))^{-1}$,
this identifies $R_0=J_\lambda*r$.  In particular
\begin{equation}\label{eq:positive-dressed-resolvent}
 R_0\ge0,\qquad
 \int_\R R_0=\frac{h}{1+h}=:\tau<\frac12.
\end{equation}
This argument uses integrable kernels, so Fourier uniqueness
identifies the actual convolution kernel, not just its multiplier.

Let $I_Q=[-Q,Q]$ and $O_Q=\R\setminus I_Q$.
The positive operator $u\mapsto R_0*(\bbone_{O_Q}u)$ has
sup norm at most $\tau$.  Its Neumann series defines bounded
nonnegative functions on the entire line by
\begin{equation}\label{eq:positive-interval-extensions}
 \begin{aligned}
 \overline r_Q&=r+R_0*(\bbone_{O_Q}\overline r_Q),\\
 \overline H_Q&=\frac1{1+h}
                     +R_0*(\bbone_{O_Q}\overline H_Q).
 \end{aligned}
\end{equation}
Multiplying by $I+K*$ and using $(I+K*)R_0*=K*$ gives
\[
 \overline r_Q+K*(\bbone_{I_Q}\overline r_Q)=s,
 \qquad
 \overline H_Q+K*(\bbone_{I_Q}\overline H_Q)=1.
\]
Their restrictions are therefore $r_Q,H_Q$, by uniqueness for
$I+K_Q$.  Since $r$ decreases on $[0,\infty)$,
the first equation in \eqref{eq:positive-interval-extensions},
restricted to $O_Q$, gives
\[
 \|\overline r_Q\|_{L^\infty(O_Q)}
 \le r(Q)+\tau\|\overline r_Q\|_{L^\infty(O_Q)},
 \qquad
 \|\overline r_Q\|_{L^\infty(O_Q)}\le\frac{r(Q)}{1-\tau}.
\]
Applying the same equation inside $I_Q$ and the second equation
on the full line now yields
\begin{equation}\label{eq:positive-interval-comparison}
 \begin{aligned}
 r(u)&\le r_Q(u)\le r(u)+\frac{\tau r(Q)}{1-\tau}
                         \le\frac{r(u)}{1-\tau},\\
 \frac1{1+h}&\le H_Q(u)\le1\qquad(|u|\le Q).
 \end{aligned}
\end{equation}
At $\lambda=\pi/2$, $K=R_0=0$, $r_Q=r=s$, and $H_Q=1$.
We have thus proved \eqref{eq:regular-density-comparison} for
every fixed $0<\lambda<\pi$.

These arguments also give the relative derivative bounds needed
for quadrature.  In the positive-kernel case, differentiate the
convolution in \eqref{eq:positive-interval-extensions}; each fixed
derivative of $R_0$ is integrable, and the exterior source is
bounded by $r(Q)/(1-\tau)$.  Thus every fixed derivative of
$r_Q-r$ is $O(r(Q))$ on $I_Q$.
For $K<0$, put $d_Q=r_Q-r$ and differentiate
$d_Q=K_{\rm out}r-K_Qd_Q$ in its output variable; the same
bound follows from $\|d_Q\|_\infty\le h r(Q)$ and the
integrability of all fixed derivatives of $K$.
Since $|r^{(j)}|\le C_jr$ and $r(u)\ge r(Q)$ on the interval,
in both cases $|r_Q^{(j)}(u)|\le C_jr_Q(u)$.
All these constants can be chosen uniformly when $\lambda$
ranges over a compact subinterval of $(0,\pi)$.
At an endpoint we extend $r_Q$ by $s-K_Qr_Q$ outside the interval.
For $K\ge0$ this is the positive function $\overline r_Q$ just
constructed; for $K<0$ it is the positive sum $s+k_Qr_Q$.
The same derivative bounds hold on an extension of bounded counting
length at either endpoint, since its physical width is
$O((Lr(Q))^{-1})$ in the dense limits used below.

Differentiation with respect to the interval endpoint, followed by
symmetry of the kernel, gives
\begin{equation}\label{eq:filling-derivative}
 \frac{\dd}{\dd Q}\int_{-Q}^Q r_Q(u)\dd u
       =2r_Q(Q)H_Q(Q)>0.
\end{equation}
Indeed $(I+K_Q)\partial_Qr_Q
=-r_Q(Q)(K(\,\cdot-Q)+K(\,\cdot+Q))$; pair this with $H_Q$.
The just proved bound $|r_Q-r|\le C r(Q)$ shows that the
integrated difference on $[-Q,Q]$ is at most $2CQr(Q)$ and
tends to zero.  Thus the mass tends to $\int_\R r=1/2$.
By \eqref{eq:regular-density-comparison}, the right side of
\eqref{eq:filling-derivative} is bounded above and below by
positive constants times $r(Q)$.  Integrating
\eqref{eq:filling-derivative} from $Q$ to infinity proves the
asserted mass-deficit estimate.

\smallskip\noindent\textbf{From dense quantiles to exact roots.}
The interval density now has a strictly increasing mass and uniform
relative derivative bounds. Choose $Q$ to have the desired mass
$n/L$, and let $y_j^0$ be the mid-quantiles of $Lr_Q$.
The continuum counting function has derivative $Lr_Q$ on the
interval and, by oddness, gives the required levels at these
quantiles.  Midpoint quadrature against a function of bounded
variation bounds the discrete residual of each root equation by
a constant.  The same holds for every fixed derivative of the
counting function: apply that quadrature to translates of the
interaction kernel or its derivatives.

Put $m_L=Lr(Q)$, and use displacement coordinates
$\zeta_j=Lr_Q(y_j^0)(y_j-y_j^0)$.  At the quantiles the Jacobian in
these coordinates is
\[
 I+\left(\frac{K(y_i^0-y_j^0)}{Lr_Q(y_j^0)}\right)_{ij}
       +O(m_L^{-1})
\]
in the sup row norm.  The middle matrix has absolute row norm
at most $h+o(1)$ by quadrature and the relative derivative bounds
just proved.  Its inverse is uniformly bounded.
On any fixed-radius ball in the $\zeta$ coordinates, the physical
displacements are $O(m_L^{-1})$.  Neighbor sums of kernel
derivatives, weighted by $1/(Lr_Q)$, are bounded.  Comparing the
bare-minus-root density to $Lr_Q$ before displacing the roots
also bounds the change of the divided diagonal entries by
$O(m_L^{-1})$.  Thus the Jacobian varies by $o(1)$ on that ball.

For clarity, Newton's map here may be taken with the inverse
Jacobian frozen at the quantiles.  Choose the ball radius larger
than twice that inverse bound times the residual bound.
For large $L$ its derivative has norm less than $1/2$, so it maps
the ball into itself and is a contraction.  The fixed point gives
the exact roots and \eqref{eq:regular-counting-density}.
Positivity of the resulting counting derivative forces their
ordering.  The same estimates hold uniformly on every fixed compact
continuation interval joining $\lambda>0$ to $\pi/2$: throughout
that interval $h$ is bounded strictly below one, and the density
and derivative comparisons above are uniform.  For the
near-half-filled sequence, the mass-deficit comparison also gives
$Lr(Q)\asymp L^\epsilon$ uniformly there.  The Jacobian is
nonsingular; equivalently its symmetrization has positive
eigenvalues, since the absolute row bound for the perturbation of
the identity is below one.

\smallskip\noindent\textbf{Quadrature and the Perron vector.}
Both root constructions are now available, with positive counting
density and a nonsingular Gaudin matrix. Their common local
quadrature estimate follows from the counting cells.
In a cell of density at least $M$, its length $\Delta_j$
is $O(M^{-1})$.  The relative derivative bound gives
$|\ell_j-\Delta_j|\le C\Delta_j^2$.  Integrate a Lipschitz test
cell by cell and include the two partial boundary cells to obtain
\eqref{eq:local-quadrature}.  Uniform local mass follows in the
dense region from the same argument and at the ends from
\eqref{eq:edge-index-density} and the gap bounds.  In regular
filling every cell is dense, including the bounded counting
extension of the interval.

The Gaudin determinant and Lemma~\ref{lem:hilbert-dual} make
each Bethe vector nonzero.  In the homogeneous chain the logarithmic
row derivative at the swap has, up to an additive scalar, local
link matrix equal to the swap times $-\eta R'(0)$.  Its hopping
coefficient on an opposite-spin pair is $-\eta/\sin\lambda<0$.
Within each nontrivial particle sector the hopping graph is
connected, so the sector minimum is simple and has a positive
vector.  At $\lambda=\pi/2$ this is the free nearest-neighbor
exclusion Hamiltonian.  In ordered particle coordinates its
eigenvectors are antisymmetrized plane waves with circle sign
$(-1)^{n-1}$, and the minimum fills the consecutive momenta about
zero.  These are precisely the consecutive Bethe roots: their
physical shift phase is
$e^{ip}=-F_{\{y\}}(i\lambda/2)$, and differentiation of the
eigenvalue gives $-2\pi\eta s(y)=-\cos p$ at the free point.
The sign gauge thus centers the occupied momenta at zero.

The real-root branch remains this simple minimum by continuation.
The physical transfer matrix is primitive in the sector: its
diagonal is positive, and successive allowed auxiliary swaps move
any particle to a neighboring hole, connecting the sector.
It commutes with the link Hamiltonian.  The positive simple
minimum is consequently its Perron vector.  Taking the long
cylinder limit gives its normalized vacuum expectations.
Finally nonsingularity of $G$ gives the finite analytic twist
branch by the implicit function theorem.  This last statement
does not require a twist neighborhood uniform in $L$.
\end{proof}


\subsection{Uniform inverses on pointwise function spaces}

Lebesgue $L^\infty$ equivalence classes cannot be used when an
operator evaluates a function at the roots.  We use $C_b(\R)$
with its pointwise sup norm, or the space of bounded continuously
differentiable functions with norm
$\|u\|_\infty+\|u'\|_\infty$.  The root-node spaces carry their
ordinary finite-dimensional sup norms.

\begin{lemma}[Gaudin and mixed-measure inverse bounds]
\label{lem:gaudin-inverse}
In the limits specified above,
\begin{equation}\label{eq:Gaudin-sup-inverse}
 \|(I+K\ell)^{-1}\|_{\infty\to\infty}\le C,\qquad
 (K\ell)_{ij}=K(y_i-y_j)\ell_j.
\end{equation}
There is also a small $a>0$ for which the inverse has uniformly
bounded row norm after conjugation by either $e^{ay_i}$ or
$e^{-ay_i}$.

Let $\chi_-,\chi_0,\chi_+$ be a nonnegative partition of unity
whose two transitions are translates of fixed smooth profiles
with bounded first derivatives.  The middle function has compact
support, the transitions lie in regions of a fixed sufficiently
high density, and the two transition regions recede from one
another.  The left function is one to the left of its transition
and zero to its right; the right function has the opposite
convention at the right transition.  Put
\[
 D_L=\nu_L-\dd u,\qquad D^\bullet=\chi_\bullet D_L.
\]
The mixed operator
\begin{equation}\label{eq:mixed-operator}
 I+R_0*(D^0+D^+)
\end{equation}
has a uniform lower bound in the pointwise sup norm.
The analogous assertion holds with the two ends interchanged.
\end{lemma}

\begin{proof}
For regular filling, local quadrature up to the endpoints bounds
the absolute row norm of $K\ell$ by $h+o(1)<1$.
This proves \eqref{eq:Gaudin-sup-inverse} there.
For $K<0$ at half filling, finite strict diagonal dominance does
not give a uniform inverse: its bare-source deficit may shrink at
the last roots. We first isolate those roots from the contracting bulk,
then use their reciprocal densities to obtain a decaying solution of
a one-sided averaging equation.

Truncate the exponentially small kernel tail at a fixed distance and
then choose a high density threshold.  Local mass and \eqref{eq:local-quadrature} show that
every row sufficiently far inside the dense portion has absolute
row sum at most a fixed $q<1$.  The edge-index bounds confine all
other rows to a fixed number of indices at either end.

If the inverse bound failed, there would be vectors $W_L$ with
$\|W_L\|_\infty=1$ and
$(I-k\ell)W_L\to0$ in sup norm.  A near-unit coordinate must
lie among those fixed edge indices.  Retain one end, translate its terminal root to zero, and extract
all fixed-index gaps and weights. At the right end write the limiting
points as $x_1=0>x_2>\cdots$; the left end is treated by reflection.  The bounds of
Lemma~\ref{lem:root-limits} make the limiting weights positive
and finite; at every fixed index they satisfy
$c/j\le\ell_j\le C/j$.  The limiting points have no finite
accumulation, since otherwise this lower bound would contradict
the uniform local mass bound.
In the array case this extraction is along $d\downarrow0$ with
$B$ sufficiently large, exactly as in
\eqref{eq:array-limit-convention}.

Put $v_j=\ell_jW_j$ in the resulting one-sided limit.
It is nonzero, and $v_j\to0$ as $j\to\infty$.
Uniform local mass and exponential tails justify passing the
equation to this limit.  It reads
\[
 (\rho_j-k(0))v_j=\sum_{i\ne j}k(x_j-x_i)v_i.
\]
The left coefficient is positive: finite row dominance bounds
it below by $k(x_j-x_i)>0$ for any other retained index $i$,
whose limiting gap is finite.  After division by this coefficient,
the equation is an
averaging equation with strictly positive off-diagonal weights
whose sum is at most one.  The latter follows from finite Gaudin
row dominance and Fatou's lemma; a strictly positive limiting
deficit is not required.  The absolute maximum of $v$ is attained,
and some other index has smaller modulus.  Its positive weight
strictly decreases this maximum, a contradiction.

We spell out the mixed case.  The measure
\[
 \widetilde\nu=\dd u+D^0+D^+
       =\chi_-\,\dd u+(1-\chi_-)\nu_L
\]
is nonnegative.  Exponential decay and bounded local variation
justify Fubini, and $R_0+K*R_0=K$ gives the exact identity
\begin{equation}\label{eq:undo-dressing}
 (I+K*)(I+R_0*(D^0+D^+))u
       =u+K*(\widetilde\nu u).
\end{equation}
The line convolution $I+K*$ is boundedly invertible because
$h<1$.

For $K<0$, suppose the undressed operator had normalized
approximate null functions $u_L$.  They satisfy
$u_L=k*(\widetilde\nu u_L)+o(1)$ uniformly.
The quantity
$\sup_x\int k(x-y)\widetilde\nu(\dd y)$ is uniformly bounded.
It follows that the supremum of $|u_L|$ on
$\supp\widetilde\nu$ is bounded below.  Divide by that supremum.
The support supremum is now one and the full supremum remains
bounded.  On the support, every row except a fixed number of
right-edge atoms has sum at most $q<1$: the replaced side is
Lebesgue, the blend is dense, and the remaining noncontracting
part is the original discrete right end.

Thus a near-unit support value survives at one of those atoms.
The modified left side is arbitrarily far away from each retained
right-edge index, so its contribution disappears in the extracted
equation.  The preceding harmonic-maximum argument applies.
There is no additional approximate mode between the atoms:
$k*(\widetilde\nu u_L)$ has a uniformly bounded derivative,
so the approximate equation makes $u_L$ equicontinuous up to a
vanishing sup error.  In the regular-filling case the entire
mixed measure instead has absolute kernel row norm $h+o(1)$,
by quadrature, and contracts directly.
Equation~\eqref{eq:undo-dressing} proves the claimed lower bound.

Finally exponential conjugation changes the discrete kernel row
norm by at most
\[
 \sup_i\sum_j |K(y_i-y_j)|\ell_j
             |e^{a(y_i-y_j)}-1|\le C|a|.
 \]
This follows by summing over unit intervals and using local mass
and exponential kernel decay.  A sufficiently small $a$ therefore
preserves the inverse bound by the inverse identity.
The same perturbation estimate preserves the mixed lower bounds.
The weighted middle contraction is proved explicitly in
Lemma~\ref{lem:weighted-flags}.
\end{proof}

\subsection{An analytic normalization for all marked rows}

The inverse estimates control the scalar-product columns.
To continue the resulting kernels from one convenient spectral
region to the whole strip, we also need a common operator bound.
Write
\[
 \mathcal S_\lambda=\{x\in\C:|\Im x|<\lambda/2\}.
\]
This is the angle strip $|\Re z|<v$ under $x=i\eta z$.

\begin{lemma}[Harmonic row normalization]\label{lem:row-normalization}
There is a real harmonic function $U$ on $\mathcal S_\lambda$
such that, for every finite real-inhomogeneity chain,
\begin{equation}\label{eq:harmonic-row-bound}
 \|T^D(x)\|\le2\|D\|
                 \exp\left(\sum_jU(x-b_j)\right).
\end{equation}
Choose a nonvanishing analytic function $\mathcal N_L$ whose
modulus is the exponential on the right.  In the homogeneous plane
and array limits, the untwisted vacuum eigenvalue $\tau_L$ satisfies
\begin{equation}\label{eq:vacuum-modulus-normalization}
 |\tau_L(x)/\mathcal N_L(x)|\longrightarrow1
\end{equation}
locally uniformly in the strip, with the moving-bulk convention
\eqref{eq:array-limit-convention} in the array case.
Consequently, for each compact $\mathcal K\subset\mathcal S_\lambda$,
\begin{equation}\label{eq:normalized-row-bound}
 \limsup\sup_{x\in\mathcal K}\|A_L^D(x)\|\le2\|D\|.
\end{equation}
In particular the vacuum denominator has no zeros on such a compact
for all sufficiently large systems in the indicated sense.
\end{lemma}

\begin{proof}
On $\Im x=\lambda/2$, local inversion says that each site matrix
divided by its equal-spin entry is unitary.  Taking the auxiliary
trace costs at most $2\|D\|$, and gives the boundary estimate with
$|a(x)|$.  Crossing gives the reflected lower boundary estimate.
For $0<\Im x<\lambda/2$ define
\begin{equation}\label{eq:harmonic-normalizer}
 U(x)=\log\left|\frac{[x+i\lambda/2]}{[i\lambda]}\right|
       +\int_\R r(u)
            \log\left|\frac{[u-x+i\lambda]}{[u-x]}\right|\dd u.
\end{equation}
The integrand has no singularities in that open half-strip and
decays exponentially at both real ends.  Its boundary value on
$\Im x=\lambda/2$ is zero, because the two sinh moduli coincide.
Reflect \eqref{eq:harmonic-normalizer} conjugately to the lower
half-strip.  The two values agree on the real line.  Differentiating
the logarithmic sinh kernels at that line, with the denominator
term interpreted by its Poisson-kernel limit, gives the upper
normal derivative
\[
 \pi(s-K*r-r)=0.
\]
Thus the two harmonic functions glue harmonically.  Their growth
at the real ends is at most linear.

For unit vectors $u,v$, apply the strip maximum principle to
$\log|\langle u,T^D(x)v\rangle|-\sum_jU(x-b_j)$.
The boundary bound is $\log(2\|D\|)$ and the growth is at most
linear.  One may justify the unbounded-strip principle by subtracting
$\varepsilon\cosh(a\Re x)\cos(a\Im x)$, with
$0<a<\pi/\lambda$, applying the principle in finite rectangles,
and then sending $\varepsilon$ to zero.
Taking the supremum over $u,v$ proves
\eqref{eq:harmonic-row-bound}.  The strip is simply connected, so
a harmonic conjugate gives the nonvanishing $\mathcal N_L$.

It remains to compare it with the actual eigenvalue.
It suffices first to work in a fixed small open neighborhood inside
the upper swap point $i\lambda/2$.
There $q_Y(x)\to0$.  Homogeneously, for each fixed
$0<\lambda<\pi$ choose a sufficiently small disc about
$i\lambda/2$.  The root ratio
$[x-y+i\lambda]/[x-y-i\lambda]$ is bounded by a constant
$M_\lambda$ for every real $y$ on that disc: the imaginary part
of its denominator stays away from $\pi\Z$, and the ratio has
finite limits at both real ends.  The one-site ratio
$[x-i\lambda/2]/[x+i\lambda/2]$ tends uniformly to zero as
the disc shrinks.  Shrink it until its supremum times
$\max(1,\sqrt{M_\lambda})$ is less than some $q_*<1$.
Then $n\le L/2$ gives $|q_Y(x)|\le q_*^L$ on the disc,
and hence on its intersection with the strict strip.
The disc and $q_*$ may depend on the fixed $\lambda$.
For the array, $\lambda>2\pi/3$ and the neighborhood can be chosen
so that each site and root ratio has modulus at most one.
Indeed the difference of the squared sinh moduli is governed,
respectively, by $\sin(2\Im x)\sin\lambda>0$ and
$\sin(2\Im x)\sin(2\lambda)<0$.
There are order $1/d$ sites at bounded distance from the moving
bulk center, with site ratios uniformly below one.  Hence
$q_Y(x)=O(e^{-c/d})$ after $B\to\infty$, uniformly in that bulk.

Put
\[
 H_x(u)=\log\left|\frac{[u-x+i\lambda]}{[u-x]}\right|,
 \qquad
 \widetilde H_x=(I+K*)^{-1}H_x.
\]
These are smooth tests in $u$ on the chosen neighborhood.
The exponentially weighted Neumann series gives uniform exponential
decay for $\widetilde H_x$ and all its fixed derivatives, also
after real recentering.  The density equation gives the exact
identity
\begin{equation}\label{eq:normalizer-discrepancy}
 \sum_{y\in Y}H_x(y)-\int_\R p_L(u)H_x(u)\dd u
       =\int_\R\widetilde H_x(u)e_L(\dd u).
\end{equation}
At half filling, integrate by parts against the bounded counting
sawtooth $F_L$.  In counting coordinates its mean over every
complete cell is zero.  On a dense fixed window, the relative
derivative bound for $\rho_L$ makes the physical cell averages
converge to that mean.  Thus $F_L$ tends weakly to zero there.
The exponential tail of the test and the uniform bound $|F_L|\le1/2$
discard the complement of an increasing window.  This proves
that \eqref{eq:normalizer-discrepancy} tends to zero, including
the array limit after bulk recentering.

For the homogeneous near-half-filled sequence, use the true
outer counting half-levels as cell endpoints.  They are a bounded
counting distance from $\pm Q$ by
\eqref{eq:regular-counting-density}, and the positive comparison
extends to them.  The cellwise primitive is zero at the endpoints
and is bounded by $1/2$ inside.  The same midpoint argument
therefore works within these cells.
Outside them, the root measure vanishes and
\[
 |\xi'(u)|\le C Lr(Q)=O(L^\epsilon).
\]
To see this uniform exterior bound, replace the root sum of any
translate of $K$ by its $Lr_Q$ integral.  Midpoint quadrature and
bounded counting displacement cost $O(1)$ by variation and
reciprocal-density summation.  The remaining exterior expression
$L(s-K_Qr_Q)$ is $O(Lr(Q))$, by the full-line equation and
\eqref{eq:finite-Q-correction}.  Its pairing with the dressed
test outside the cells is consequently
$O(L^\epsilon e^{-aQ})$ for a fixed $a>0$.
Since $Q=\eta(1-\epsilon)\log L+O(1)$, it tends to zero when
\begin{equation}\label{eq:epsilon-choice}
 0<\epsilon<\frac{a\eta}{1+a\eta}.
\end{equation}
Decreasing $a$ if necessary gives one such choice on the compact
continuation interval used to construct the roots.

The first-term eigenvalue formula, $q_Y\to0$, and
\eqref{eq:normalizer-discrepancy} now prove
\eqref{eq:vacuum-modulus-normalization} on this open neighborhood.
On the whole strip, $\tau_L/\mathcal N_L$ is bounded by two.
Every subsequential holomorphic limit therefore has modulus one
on the neighborhood and is a unimodular constant by the maximum
modulus principle.  Normal-family compactness gives local uniform
modulus convergence throughout the strip.  The same argument
works for any purported failing moving-bulk array sequence.
Dividing \eqref{eq:harmonic-row-bound} by $|\tau_L|$ proves the
last assertions.
\end{proof}

\subsection{Joint tails from finite residue conditions}

Fix a word length $m$ and generic arguments $v'_1,\ldots,v'_m$
in a small complex neighborhood of $i\lambda/2$.  Generic here
means that the finite difference factors in
Corollary~\ref{cor:finite-word-expansion} have no poles or forbidden
coincidences.  This is a nonempty open condition.
We may recenter all real coordinates in the array case.
The columns in \eqref{eq:W-determinant}, at zero twist, are denoted
by $W_{v'}$.  Their source columns are
\begin{equation}\label{eq:thermodynamic-source-column}
 m_{v'}(u)=
 \frac{(1+q_Y(v'))\coth(v'-u)-\coth(v'-u-i\lambda)
                       -q_Y(v')\coth(v'-u+i\lambda)}
      {-2\pi i}.
\end{equation}
Extend each finite node solution to the line by
$W=m-K*(\nu_LW)$.  Lemma~\ref{lem:gaudin-inverse}, local mass,
and the kernel derivative bounds give a uniform bound for $W$
and each fixed derivative on the line.  Undoing this convolution
equation gives
\begin{equation}\label{eq:finite-dressing-equation}
 W=W^0-R_0*(D_LW),\qquad W^0=(I+K*)^{-1}m.
\end{equation}
Write $E=\C^m$ for the coefficient space and
$W(c)=\sum_{j=1}^m c_jW_{v'_j}$ for $c\in E$, with the same
notation for linear combinations of the free columns $W^0_{v'_j}$.
All these statements are uniform for $\|c\|_2\le1$.

We need decay for a whole minor of these columns.  To see how much,
fix a finite action path and write $H=(p_1,\ldots,p_k)$ for its deleted
original roots and $V$ for its surviving row arguments.  In
\eqref{eq:exact-path-determinant}, leave the determinant
$\det(W_{v'}(p))_{H,V}$ aside for the moment and extract the
quadrature weights $\prod_{p\in H}\ell_p$.
The hyperbolic Cauchy determinant in
\eqref{eq:exact-path-determinant} contributes at most
\[
 C\exp\left(k\sum_i|p_i|-\sum_{i<j}|p_i-p_j|\right)
\]
after the root-pair cancellations.  To check those cancellations
also near collisions, for each pair of deleted original roots
Corollary~\ref{cor:finite-word-expansion} supplies a reciprocal
$f$ factor.  Combined with its Cauchy denominator, its modulus
is bounded by
\[
 |[p_i-p_j\pm i\lambda]|^{-1}
       \le C_\lambda e^{-|p_i-p_j|}.
\]
All other fixed-argument factors are bounded on the generic
parameter set.  There is at least one removal factor
$\beta(v'-p_i)$ or $\beta(p_i-v')$ for each original root
deleted, and it is at most $Ce^{-|p_i|}$.
The factors $q_Y(v')$ are bounded and the row denominators
$1+q_Y(v')$ are bounded away from zero in the spectral region
used here.  Consequently, apart from the dressed determinant and the quadrature
weights, the summand has modulus at most
\[
 C\exp\left((k-1)\sum_i|p_i|
                         -\sum_{i<j}|p_i-p_j|\right).
\]
With $k_+(t)=\#\{i:p_i>t\}$ and
$k_-(t)=\#\{i:-p_i>t\}$, its exponent is exactly
\begin{equation}\label{eq:Cauchy-tail-cost}
 \int_0^\infty
 \bigl(k_+(t)(k_+(t)-1)+k_-(t)(k_-(t)-1)\bigr)\dd t.
\end{equation}
One verifies this identity first on a region with the signs and
ordering fixed, or by writing every distance as an integral of
indicator functions.

Thus $r$ roots escaping together at one end cost $r(r-1)$ in the
exponential rate.  A bound on each column by the same decaying function
would not control arbitrarily large fixed minors.  We will instead find
nested subspaces of column combinations with successively faster decay.
For the $j$th direction it suffices to obtain a rate
$\beta_j>2j-2$, because the first $r$ such rates then sum to more than
$r(r-1)$.  The number of slow directions is determined by the poles of
the full-line dressed equation.

The pole spaces needed for this equation are particularly small.
Let $(d_j)$ be the increasing list, with multiplicities, formed
from
\begin{equation}\label{eq:decay-pole-list}
 \left\{\frac{2k-1}{\eta}:k\ge1\right\}
 \quad\cup\quad
 \left\{\frac{2k}{1-\eta}:k\ge1\right\}.
\end{equation}
At coincident entries, retain both copies; the list bounds the
maximum possible pole orders.  For every $j$,
\begin{equation}\label{eq:pole-count}
 d_j\ge2j-1.
\end{equation}
Indeed, if $n_1,n_2$ count the entries of the two progressions
strictly below $T=2j-1$, then
$2n_1-1<\eta T$ and $2n_2<(1-\eta)T$.
Adding gives $n_1+n_2<j$.

For completeness, the Fourier calculation behind the list uses
\[
 1+\widehat K(p)
 =\frac{2\sinh((\pi-\lambda)p/2)\cosh(\lambda p/2)}
        {\sinh(\pi p/2)}.
\]
In particular
\[
 \widehat R_0(p)=
 \frac{\sinh((\pi-2\lambda)p/2)}
      {2\sinh((\pi-\lambda)p/2)\cosh(\lambda p/2)}.
\]
The Fourier transform of each coth difference in
\eqref{eq:thermodynamic-source-column} has denominator
$\sinh(\pi p/2)$, as follows by shifting the integration contour
by $i\pi$ and accounting for the coth residues.
This denominator cancels upon dressing.  The remaining possible
poles for both $W^0$ and $R_0$ are exactly those bounded by
\eqref{eq:decay-pole-list}; the apparent singularity at zero is
removable.  On any horizontal Fourier contour avoiding these
poles, the transforms decay exponentially in the real direction.
The generic arguments range in a compact analytic neighborhood,
so this decay is uniform.  Multiplication by any fixed power of
the Fourier variable preserves integrability.

Thus, for $\beta>0$ not in the list, shifting the contour and
subtracting the crossed residues gives
\[
 R_{0,\beta}=R_0-P_\beta,\qquad
 W^0_\beta(c)=W^0(c)-p_\beta(c),
\]
where the global exponential polynomials $P_\beta,p_\beta(c)$
belong to
\begin{equation}\label{eq:residue-space}
 \mathcal E_\beta
 =\Span\{u^q e^{-du}:d<\beta,\ 0\le q<m_d\}.
\end{equation}
Here $m_d$ is the multiplicity in the list, and
$\dim\mathcal E_\beta\le r_\beta:=\#\{j:d_j<\beta\}$.
Using two nearby pole-free contours gives, for some $\delta>0$,
\begin{equation}\label{eq:shifted-kernel-envelope}
 e^{\beta u}|R_{0,\beta}^{(q)}(u)|
       \le C_{\beta,q}e^{-\delta|u|}
 \quad(q=0,1,2),
\end{equation}
and the corresponding bounds for the free columns.
Any other fixed derivative order is available in the same way.
The residue polynomials are subtracted on the whole line, without
cutoffs; this preserves their finite-dimensional translation span.

\begin{lemma}[Finite-chain weighted flags]\label{lem:weighted-flags}
Let $E=\C^m$ be the coefficient space of the fixed column list.
For each admissible $\beta>0$ there is a subspace
$F_{\beta,L}\subset E$ of codimension at most $r_\beta$
such that
\begin{equation}\label{eq:weighted-node-decay}
 |W(c)(u)|\le C_\beta e^{-\beta u}\|c\|_2
 \quad\hbox{at every positive root }u,\quad c\in F_{\beta,L}.
\end{equation}
For a fixed finite list of increasing shifts, choose one discrepancy
partition for the whole list. The resulting subspaces may be chosen
nested.  There are analogous flags at the negative end.
The constants are uniform in the three limiting procedures,
with regular filling understood as $Q\to\infty$ after $L\to\infty$.
Restriction to any column coordinate subspace preserves the
codimension bound and all unit-vector estimates.
\end{lemma}

\begin{proof}
\smallskip\noindent\textbf{Separate the two sparse ends.}
Choose the fixed-profile partition in Lemma~\ref{lem:gaudin-inverse}
so that
\[
 \supp D^-\subset(-\infty,a_L],\qquad
 \supp D^+\subset[b_L,\infty),\qquad a_L<0<b_L.
 \]
The word center is deep between these boundaries.  The middle
support is compact and dense.  At a fixed density threshold,
Lemma~\ref{lem:root-limits} gives
\begin{equation}\label{eq:last-node-split}
 \max Y\le b_L+C_0,\qquad \min Y\ge a_L-C_0,
\end{equation}
with a threshold-dependent constant uniform in the limits.
For regular filling take the transitions just inside $[-Q,Q]$.
For the array both the last root and the transition contain the
same $\eta\log(1/d)$ shift in \eqref{eq:nominal-edges}.
All three signed pieces have uniformly bounded variation mass
on intervals of length one.

Use the pointwise Banach space
\[
 X_\beta=\left\{u\in C^1(\R):
       \|u\|_\beta:=\sup_x e^{\beta x}
                  (|u(x)|+|u'(x)|)<\infty\right\}.
\]
The quadrature error on each unit-scale test is as small as desired
on the middle support and its taper.  To bound
$R_{0,\beta}*(D^0u)$, insert a smooth unit-scale partition in the
integration variable and apply \eqref{eq:local-quadrature}.
For the output and its first derivative, the test derivatives
involve $R_{0,\beta},R_{0,\beta}',R_{0,\beta}''$ and $u,u'$.
The envelopes \eqref{eq:shifted-kernel-envelope} sum after
weighting to give
\[
 \|R_{0,\beta}*(D^0u)\|_\beta
       \le C_\beta\varepsilon_{\rm quad}\|u\|_\beta.
\]
Choose the density threshold so this norm is below $1/2$ for the
fixed finite list of shifts, including zero.

We first remove the right-end discrepancy from the column equation,
keeping its true left-end input.  Let
\[
 W^p(c)=(I+R_0*D^0)^{-1}
             (W^0(c)-R_0*(D^-W(c))),
\]
using the shift-zero contraction on $X_0$.
Our task is to characterize the combinations $c$ for which this
bounded column also belongs to $X_\beta$.  The omitted right end will
then be restored with the mixed-operator lower bound.

\smallskip\noindent\textbf{Solve after subtracting the slow residues.}
For $c\in E$ solve in $X_\beta$
\begin{equation}\label{eq:shifted-middle-equation}
 (I+R_{0,\beta}*D^0)U_\beta(c)
     =W^0_\beta(c)-R_{0,\beta}*(D^-W(c)).
\end{equation}
The true-column bound and local variation imply
\[
 \|R_{0,\beta}*(D^-W(c))\|_\beta
       \le C_\beta e^{\beta a_L}\|c\|_2.
\]
For example multiply the convolution integrand by
$e^{\beta(x-y)}e^{\beta y}$, use $y\le a_L$, and sum the
exponentially decaying weighted kernel envelopes over unit
intervals.  The same argument applies to the first derivative.
Consequently
\begin{equation}\label{eq:shifted-solution-bound}
 \|U_\beta(c)\|_\beta\le C_\beta\|c\|_2.
\end{equation}
This estimate holds for every coefficient vector before imposing
any residue condition.

\smallskip\noindent\textbf{Impose only the residue conditions.}
Define the linear map
\begin{equation}\label{eq:residue-condition-map}
 \Lambda_\beta(c)
  =P_\beta*(D^0U_\beta(c)+D^-W(c))-p_\beta(c).
\end{equation}
Its values lie in $\mathcal E_\beta$.  Indeed convolution of an
exponential polynomial with a source preserves that span:
expand each $(x-y)^q$.  The middle source is compactly supported,
and the left moments converge because $d>0$ and the left
variation mass is locally bounded.  Rank-nullity therefore gives
$\operatorname{codim}\ker\Lambda_\beta\le r_\beta$.
No inverse or condition-number bound for $\Lambda_\beta$ is used.

Restoring the residues in \eqref{eq:shifted-middle-equation} gives
\[
 (I+R_0*D^0)U_\beta(c)
  =W^0(c)-R_0*(D^-W(c))+\Lambda_\beta(c).
\]
For $c\in\ker\Lambda_\beta$, this identity first proves
$U_\beta(c)\in X_0$: in the rearranged equation the middle
source $D^0U_\beta$ is a finite compactly supported measure,
so its convolution and first derivative are bounded at both
ends.  The other two terms are already bounded.
This membership assertion for a fixed chain does not yet require
a uniform estimate for that compact source.  Only now apply the
shift-zero inverse to identify $U_\beta(c)=W^p(c)$ and obtain
its uniform unweighted bound.

Conversely, if $W^p(c)\in X_\beta$, rewrite its equation with
the shifted kernels.  The defect is an element of
$\mathcal E_\beta$ which also belongs to $X_\beta$.
It is zero: every nonzero exponential polynomial in
\eqref{eq:residue-space} violates the $\beta$-weighted bound at
$+\infty$.  Weighted uniqueness then gives
$W^p(c)=U_\beta(c)$ and $\Lambda_\beta(c)=0$.
Hence
\begin{equation}\label{eq:flag-membership}
 F_{\beta,L}:=\ker\Lambda_\beta
       =\{c:W^p(c)\in X_\beta\}.
\end{equation}
The spaces are nested, because $W^p(c)$ is already unweighted
bounded at $-\infty$ and a larger positive weight at $+\infty$
implies every smaller one.  Estimates
\eqref{eq:shifted-solution-bound} hold uniformly for every unit
vector in these subspaces.

\smallskip\noindent\textbf{Restore the right end at the actual nodes.}
For such a vector, subtract the true and middle equations:
\[
 (I+R_0*(D^0+D^+))(W(c)-W^p(c))
       =-R_0*(D^+W^p(c)).
\]
The right side has sup norm at most $C_\beta e^{-\beta b_L}$.
The infinite Lebesgue portion of $D^+$ causes no extra factor:
sum its uniformly bounded local variation against
$e^{-\beta(b_L+j)}$ over $j\ge0$.
The mixed lower bound therefore gives the same estimate for
$W-W^p$.  At every actual positive root,
\eqref{eq:last-node-split} converts it to
$C_\beta e^{-\beta u}$, which proves
\eqref{eq:weighted-node-decay}.
Reflecting proves the negative-end assertion.
Finally restriction of \eqref{eq:residue-condition-map} to a
coordinate subspace still has rank at most $r_\beta$, and a unit
vector there is a unit vector in $E$.  All the asserted
restriction properties follow.
\end{proof}

The flag estimates now supply the decay rates required by
\eqref{eq:Cauchy-tail-cost}.

\begin{lemma}[Joint tail bound]\label{lem:joint-determinant-tails}
For a fixed generic row word and any one of its nonzero finite
action paths, let $H=(p_1,\ldots,p_k)$ be the deleted original
roots and $V$ its surviving row arguments.
After extracting the quadrature weights $\prod_{i=1}^k\ell_{p_i}$
from \eqref{eq:exact-path-determinant}, the remaining summand is
bounded in modulus by
\begin{equation}\label{eq:joint-tail-majorant}
 C\exp\left(-\varepsilon\sum_{i=1}^k|p_i|\right)
\end{equation}
for some $\varepsilon>0$.  The constant and exponent may depend
on the fixed word and generic compact parameter set, but are
uniform in the limiting procedures.  In particular, the sum over
paths with $\max_i|p_i|>R$ is $O(e^{-\varepsilon'R})$ uniformly,
for some $\varepsilon'>0$.
\end{lemma}

\begin{proof}
Choose pole-free numbers
$\beta_j\in(2j-2,2j-1)$ for $1\le j\le k$.
By \eqref{eq:pole-count}, the corresponding flag has
codimension at most $j-1$.
For a positive-row minor with $r$ columns, order its rows as
$p_1\ge\cdots\ge p_r\ge0$.
Choose an orthonormal basis of its coefficient space backwards:
the $j$th basis vector belongs to $F_{\beta_j,L}$.
This is possible because these spaces are nested and have
dimension at least $r-j+1$.  A unitary basis change does not alter
the determinant modulus.  Applying
\eqref{eq:weighted-node-decay} to its columns and using the
rearrangement inequality in the determinant expansion gives
\begin{equation}\label{eq:positive-minor-bound}
 |\det W|\le C\exp\left(-\sum_{j=1}^r\beta_jp_j\right)
   =C\exp\left(-\int_0^\infty
                \sum_{j=1}^{k_+(t)}\beta_j\dd t\right),
\end{equation}
where $k_+(t)=\#\{i:p_i>t\}$.
For a determinant with both signs, first expand by its positive
and negative rows.  Every Laplace minor has a column coordinate
subspace, on which the same codimension and unit-vector estimates
hold.  Choose its flag basis separately.  Thus no simultaneous
adaptation of the two tail flags is required.  The finite Laplace
sum gives the product of the positive bound and its negative
counterpart.

Set $\varepsilon=\min_{j\le k}(\beta_j-(2j-2))>0$.
For every $r\le k$,
$\sum_{j=1}^r\beta_j\ge r(r-1)+\varepsilon r$.
Combining \eqref{eq:positive-minor-bound} at both ends with
\eqref{eq:Cauchy-tail-cost} proves
\eqref{eq:joint-tail-majorant}.  The case $k=0$ is the empty
determinant and needs no tail estimate.
Uniform local quadrature mass now bounds the sums of the
majorant by convergent geometric series on unit boxes.
Discarding boxes with some coordinate beyond $R$ proves the
last assertion, with a harmless smaller exponent.
\end{proof}

\subsection{Limits of fixed row words}

The joint tail bound lets us take the thermodynamic limit first on a
bounded rapidity window, and then remove that window.  The harmonic
row bound extends the resulting convergence to the whole strict strip.

\begin{proposition}[Limits of fixed row words]\label{prop:word-kernel}
For each fixed integer $m\ge0$ and fixed $2\times2$ matrices
$D_1,\ldots,D_m$, the homogeneous plane coefficients
\begin{equation}\label{eq:word-kernel-definition}
 \left\langle\Omega_L,
     A_L^{D_m}(x_m)\cdots A_L^{D_1}(x_1)\Omega_L\right\rangle
\end{equation}
converge locally uniformly on $\mathcal S_\lambda^m$ to a
holomorphic function
$\mathcal K_{D_m,\ldots,D_1}(x_m,\ldots,x_1)$.
It is invariant under every common complex translation for which
the translated arguments remain in the strip.
For $\lambda>2\pi/3$, the same functions are the array limits in
\eqref{eq:array-limit-convention}.

If $v_L$ is any fixed finite linear combination of row words applied to
$\Omega_L$, then
\[
 \lim\|A_L^D(x)v_L\|^2
       \le4\|D\|^2\lim\|v_L\|^2.
\]
The two limits are determined by the word coefficients and the adjoint
rule. This is the bound used to construct bounded row actions in
Section~\ref{sec:angular}.
For $\lambda>2\pi/3$, in the regular-filling procedure the expectations of every fixed
equal-time local matrix product converge to those of the
homogeneous plane construction.
\end{proposition}

\begin{proof}
First keep all arguments in the generic neighborhood used above.
The finite row expansion has finitely many action patterns, each
with at most $m$ deleted original roots.  Words with a net change
of particle number have coefficient zero.
On each bounded rapidity window, the line extensions of the
dressed columns and their derivatives are uniformly bounded.
Local quadrature and their convolution equation therefore give,
along any subsequence, a smooth local limit satisfying
\[
 W+K*W=m^{(0)}_{v'},
 \qquad
 m^{(0)}_{v'}(u)=
       \frac{\coth(v'-u)-\coth(v'-u-i\lambda)}{-2\pi i}.
\]
The tails in this equation are negligible by exponential kernel
decay and bounded local mass.  Since $h<1$, the bounded solution
is unique.  The elementary identity
$m^{(0)}_{v'}(u)=s(u-v'+i\lambda/2)$ shows that dressing
replaces $s$ by $r$.  Equivalently its Fourier transform is
$e^{-ip(v'-i\lambda/2)}/(2\cosh(\lambda p/2))$.
This gives
\begin{equation}\label{eq:continuum-dressed-column}
 W_{v'}(u)\longrightarrow
 \frac{1}{2\lambda
          \cosh((u-v'+i\lambda/2)/\eta)}.
\end{equation}
For the array, the same conclusion holds along every bulk-recentered
sequence in its prescribed double-limit sense.

Apply \eqref{eq:exact-path-determinant} on a bounded rapidity
window.  Its weights are exactly the product quadrature
measure.  Local quadrature, \eqref{eq:continuum-dressed-column},
and $q_Y(v')\to0$ turn its sum into the corresponding line
integral.  The apparent near-collision singularities have already
been canceled by the root-pair factors; deleting diagonal tuples
has no effect in the limit, since the maximal atom on such a
window tends to zero.  Lemma~\ref{lem:joint-determinant-tails}
then permits the window to grow to the whole line.
This proves convergence of every path and hence of the word.
It also proves independence of the choice between homogeneous
and array limits.  All surviving ingredients depend only on
rapidity differences.  Changing all variables of integration
by a common real amount therefore proves real-translation
invariance in the $x$ coordinate.

Lemma~\ref{lem:row-normalization} bounds all the finite words
locally uniformly on the full strip.  Normal-family compactness
and the identity theorem extend the convergence from the generic
open set to $\mathcal S_\lambda^m$, with uniqueness of every
subsequential limit.  The derivative in the common translation
direction vanishes on that open set and hence everywhere.
Integrating it along translations staying in the strip gives the
stated complex-translation invariance.
For any finite linear combination of word vectors, apply the
finite operator inequality \eqref{eq:normalized-row-bound} to
its Gram norm.  Passing to the limit gives the bound
$2\|D\|$ for the induced row action and also shows that it
respects zero Gram vectors.

For the remaining equal-time assertion, first fix $Q<\infty$
and send $L\to\infty$ with filling
$\int_{-Q}^Qr_Q$.  The same finite path calculation is valid,
now with dense nodes on $[-Q,Q]$ and the interval dressing
operator.  Then send $Q\to\infty$.
The inverse bounds, local convergence, and joint tail bound
are uniform in this second limit, so the same full-line integrals
result.

This calculation can be made on small contours surrounding
$i\lambda/2$, using nested distinct radii to keep all arguments
generic.  For the regular-filling assertion we only use
$\lambda>2\pi/3$; choose the contour radii so small that all the
fixed shifted poles stay outside them.
On these contours the homogeneous $q_Y$ still tends to zero
exponentially, and the eigenvalue first term is analytic and
nonzero.  Thus each complete finite normalized word is holomorphic on the
product of the enclosed discs for sufficiently large $L$.
The same determinant bounds apply on their boundary contours, including
the portions just outside the strict strip. Convergence on these
contours is uniform: first take the fixed-$Q$ chain limit on a bounded
rapidity window, then use the uniform joint tail bound when taking
$Q\to\infty$ and removing that window. The iterated Cauchy formula
therefore gives the coincident endpoint values.

At the swap, a marked row implements an arbitrary matrix on one site
together with the cyclic shift. Multiplying marked swap rows conjugates
the successive site matrices by that shift; their product reads any
finite set of sites, with identity marks at unused sites. The residual
shift fixes the homogeneous Perron vector, so its vacuum expectation
is exactly the desired local-matrix expectation. The half-filled
homogeneous calculation has the same endpoint contours and integrals.
This proves
equality of the equal-time local states in the stated order of
limits, without asserting an unrestricted row-strip limit at a
fixed nonzero magnetization.
\end{proof}

\subsection{Clustering of diagonal row words}

\begin{proposition}[Separated diagonal words]\label{prop:word-clustering}
Partition a fixed finite list of diagonal marked rows into groups.
Within each group fix its internal arguments in the strict angle
strip.  Translate the arguments of group $g$ by a common imaginary
amount $it_g$, and let every $|t_g-t_{g'}|$ tend to infinity.
Then the plane word kernel factors into the product of the
individual group kernels.  Convergence is locally uniform in
the internal strict arguments.  The order within each group is
retained.
\end{proposition}

\begin{proof}
We first use generic arguments near the upper boundary in the
$x$ coordinate.  Their real parts may be separated by arbitrary
translations.  The integral formula of
Proposition~\ref{prop:word-kernel} applies there by analytic
continuation of its difference integrals; keep the imaginary
parts in a fixed small generic box and avoid the finitely many
coincidence poles during continuation.

For a path with deleted roots $H$ and surviving arguments $V$,
write $N_H(t)$ and $N_V(t)$ for the numbers of their real parts
to the left of $t$.  These two lists have equal length.
The quotient of the continuum determinant
\eqref{eq:continuum-dressed-column} and the original hyperbolic
Cauchy determinant, including the root-pair cancellations, is
bounded by
\begin{equation}\label{eq:matching-cluster-bound}
 C\exp\left(-(\eta^{-1}-1)
                   \int_\R(N_H(t)-N_V(t))^2\dd t\right).
\end{equation}
Here the constant is uniform as the fixed internal lists separate.
To verify the exponent, use
\[
 \cosh((p-v'+i\lambda/2)/\eta)
       =i\sinh((p-v')/\eta)
\]
and apply the sinh Cauchy identity at scales $1/\eta$ and $1$.
The exponent in their ratio is $-(\eta^{-1}-1)$ times
\[
 \sum_{p\in H,v'\in V}|p-\Re v'|
       -\sum_{p<p'}|p-p'|
       -\sum_{v'<v''}|\Re v'-\Re v''|.
\]
This expression is exactly the integral in
\eqref{eq:matching-cluster-bound}.  The bounded imaginary offsets
and fixed within-group generic factors change only the constant.
The reciprocal $f$ factors remove the near-root-collision
denominators as in Lemma~\ref{lem:joint-determinant-tails}.

The integral in \eqref{eq:matching-cluster-bound} controls ordered
matching distance.  Indeed $N_H-N_V$ is integer-valued, so its
square is at least its absolute value, and
\[
 \int_\R|N_H-N_V|\dd t
     =\sum_j|p_j-\Re v'_j|
\]
for the two ordered lists.  Since $\eta<1$, the estimate
localizes all integration variables integrably around the
surviving arguments, uniformly in their separations.

For diagonal words every action has one removal.  Each removal
also has its factor $\beta$ decaying exponentially in the
distance between the adding argument and the removed argument
or root.  If a path couples two different groups, either this
factor or \eqref{eq:matching-cluster-bound} makes its integral
vanish as the gap grows.  More explicitly, first restrict every
matched root to a fixed radius of its matching argument using
the uniform exponential bound.  On this restricted region any
removal between distinct groups has distance tending to infinity.
Discard these paths and then increase the matching radius.
The surviving paths have their additions and removals inside
individual groups; in each group the numbers of deleted original
roots and surviving added arguments are equal.

On a bounded matching region the two determinants are block
diagonal in the limit, so their quotient tends to the product of
the group quotients.  Cross-group correction products tend to
one.  In fact $f(u)\to e^{\pm i\lambda}$ as
$u\to\pm\infty$, and at every removal the numbers of extra
surviving arguments and missing original roots in another group
agree.  Their limiting phases cancel.  The constants
$(-[i\lambda])^{-1}$ occur once per deleted column and multiply
in the same way.  Termwise limits, the integrable matching bound,
and the finite sum over paths therefore give exactly the product
of the separate group answers.

If different groups were interlaced in the original word, they
may first be gathered while retaining each internal order.
The normalized intertwiner in RTT tends exponentially to an
invertible diagonal matrix at large real separation.
Its conjugation of a pair of diagonal marks therefore differs
from the unchanged pair by an exponentially small matrix.
The bounded limiting rows make the corresponding exchange error
tend to zero.

Finally the row norm bound makes these translated kernels a
normal family as functions of the internal strict arguments.
Every subsequential holomorphic limit equals the product on the
generic open box just treated.  The identity theorem determines
it throughout the product strip, proving the proposition and
its local uniformity.
\end{proof}

The construction now supplies all row-kernel inputs for the
angular Hilbert space: holomorphy, the common-translation
symmetry, bounded marked actions, and diagonal clustering.
The physical identification and endpoint continuation are
treated next.
